\begin{abstract}
Evaluation games read the quantifiers of first-order logic as moves by two players: one chooses witnesses for $\exists$, the other for $\forall$. What do such games mean with more than two players, when each choice belongs to one of them and each play has one winner? We take the value of a finite game tree to be the set of players who win under rational play, that is, in some pure subgame-perfect equilibrium. Values compose by a one-line rule, giving a finite algebra of owned choices that is Boolean for two players and not for three.

We determine much of its structure. An equation holds as soon as it holds when each variable stands for at most two winners, and this bound is optimal. The idempotent operations are exactly those built from owned choices alone. Deciding equations is coNP-complete, and the boundary of tractability is sharp: owned choice, union, and a derived operation called reset are each easy alone and hard in any pair, while reset with one player's guarded choice stays polynomial. If moves may be hidden until the end of a round, values must become sets of probability distributions, and no finite compositional semantics can track even who might win. Key results are formalized in Lean.
\end{abstract}
